\documentclass[10pt,a4paper]{amsart}
\usepackage[margin=19mm]{geometry}
\usepackage{amsmath,amssymb,mathtools}
\usepackage[scr=rsfs,cal=euler]{mathalfa}
\usepackage{xfrac}
\usepackage[unicode,hidelinks,bookmarks=false]{hyperref}
\newcommand{\ie}{i.e.\ }
\newcommand{\cf}{cf.\ }
\newcommand{\resp}{respectively\ }
\newcommand{\Subsection}{\subsection}
\providecommand{\nobreakdash}{\nobreak\hbox{-}}
\setlength{\emergencystretch}{2em}
\multlinegap0pt
\usepackage{bm}
\usepackage{bbm}
\usepackage{dsfont}
\usepackage{xspace}
\usepackage{cancel}
\usepackage{mathtools}
\usepackage{amsthm}
\makeatletter
\@ifundefined{thm}{\newtheorem{thm}{Thm}}{}
\@ifundefined{lem}{\newtheorem{lem}[thm]{Lemma}}{}
\makeatother
\def\d{\,\mathrm{d}}
\title[Sharp asymptotic behaviour of symmetric and non-symmetric so]{Sharp asymptotic behaviour of symmetric and non-symmetric solutions of the Heat Equation in the Hyperbolic Space}
\author{Jos\'e Alfredo Ca\~nizo, Alejandro G\'arriz, Diego Alfonso Mar\'in}
\date{Source: arXiv:2604.13669v1 (CC BY 4.0)}
\begin{document}
\maketitle

\noindent\emph{Source and licence.} Sharp asymptotic behaviour of symmetric and non-symmetric solutions of the Heat Equation in the Hyperbolic Space. Version arXiv:2604.13669v1 (CC BY 4.0), distributed under the Creative Commons Attribution 4.0 International licence, as recorded in the arXiv licence metadata for this exact version and shown on the abstract page https://arxiv.org/abs/2604.13669v1. Full rights notice: licenses/arxiv-2604.13669-CC-BY-4.0.txt.\par\smallskip

\noindent\textit{Editorial scope.} Counted unit: the Lem whose source label is \texttt{lem:constant\_C(t)} in the article (source0.tex lines 2287--2304, source label \texttt{lem:constant\_C(t)}), delivered together with the article's own complete original proof of it (source0.tex lines 2305--2437, one \texttt{proof} environment of 133 lines). No mathematics is written, repaired or re-derived: statement and proof are the article's own text line for line. Required context retained uncounted: eq:definition\_bottom\_spectrum\_laplacian (lines 272--274); eq:main\_spherical\_coordinates\_symmetric\_datum (lines 656--662); eq:bound\_quotient\_C(t) (lines 682--684); eq:limit\_C(t) (lines 690--692). Every definition, display and label of those lines is retained unchanged. Omitted from this package and not redistributed: 1--271, 275--655, 663--681, 685--689, 693--2286, 2438--2842. Nothing inside a retained excerpt is omitted or altered: every retained body is delivered exactly as the article's lines are written, so no omission receipt of any kind accompanies this package. Citations: the retained excerpts carry no citation command at all, so no bibliography entry of the article is transcribed and no reference is redistributed. Reference adaptation: none was needed and none was made; every reference inside the delivered unit resolves inside it, so no unresolved reference or citation is produced. Printed slips kept literal: no printed word, symbol, hypothesis or number of the counted statement or of its proof was altered. Layout and class: the article's own class is \texttt{article} with packages including inputenc, fontenc, geometry, amsmath, amsthm, amsfonts, amssymb, mathrsfs, parskip, hyperref, cleveref, bbm, natbib, url; the wrapper is the lane's pinned \texttt{amsart} editorial wrapper at 10pt on A4, which restates the theorem-like environments the retained text needs and adds the article's own macro definitions that the retained text needs (\texttt{\textbackslash d}), with extra wrapper packages (bm, bbm, dsfont, xspace, cancel, mathtools). The delivered numbering is the unit's own. Assets: no retained span contains a figure environment, a graphics-inclusion command, a PDF-page inclusion, a table or a TikZ picture; no image file is copied into this package, the package declares no asset dependency and no image file is redistributed with it. Licence: version arXiv:2604.13669v1 of this article is distributed under the Creative Commons Attribution 4.0 International licence, as recorded in the arXiv licence metadata for this exact version and shown on the abstract page https://arxiv.org/abs/2604.13669v1. A full rights notice is delivered at licenses/arxiv-2604.13669-CC-BY-4.0.txt. Characters: every retained excerpt is pure ASCII, so the delivered engine, XeLaTeX with its default Latin Modern text font, reports no missing character.
\begin{equation}\label{eq:definition_bottom_spectrum_laplacian}
	\lambda_1:=\frac{(d-1)^2}{4}.
\end{equation}

\begin{equation}\label{eq:main_spherical_coordinates_symmetric_datum}
\begin{cases}
	v_\tau = \partial_\rho\left(v_\rho + \left[(d-1)e^{\frac{\tau}{2}}+\frac{1}{2}\rho- (d-1)e^{\frac{\tau}{2}}\coth(r)\right]  v\right),\quad &\rho\geq \rho_0(\tau),\ \tau>0,\\
	v(\tau, \rho_0(\tau))=0, &\tau>0,\\
	v(0,\cdot) = v_0, &\rho\geq \rho_0(\tau).
\end{cases}
\end{equation}

	\begin{equation}\label{eq:bound_quotient_C(t)}
	-k_d\  e^{-m_d (t+1)}\leq \frac{\partial_t C(t)}{C(t)}\leq k_d\  e^{-m_d (t+1)},
\end{equation}

\begin{equation}\label{eq:limit_C(t)}
	\frac{2^{d-2}\mathcal{M}_{\mathcal{R}}}{\sqrt{\pi} }\leq C(t)\leq \frac{2^{d-2}\mathcal{M}_{\mathcal{R}}}{(1-e^{-(d-1)(t+1)})^{d-1}\left[\sqrt{\pi} - \frac{1}{\sqrt{\lambda_1 (t+1)}}e^{-\frac{\lambda_1 (t+1)}{4}}\right]}.
\end{equation}

\begin{lem}\label{lem:constant_C(t)}
	There exists a constant $k_d$ depending only on the dimension $d$ such that
	\begin{equation*}
		-k_d\  e^{-m_d (t+1)}\leq \frac{\partial_t C(t)}{C(t)}\leq k_d\  e^{-m_d (t+1)},
	\end{equation*}
	where
	\[
	m_d := \min \left(d-1, \frac{(d-1)^2}{16}\right).
	\]
	Moreover,
	\begin{equation*}
		\frac{2^{d-2}\mathcal{M}_{\mathcal{R}}}{\sqrt{\pi} }\leq C(t)\leq \frac{2^{d-2}\mathcal{M}_{\mathcal{R}}}{(1-e^{-(d-1)(t+1)})^{d-1}\left[\sqrt{\pi} - \frac{1}{\sqrt{\lambda_1 (t+1)}}e^{-\frac{\lambda_1 (t+1)}{4}}\right]}
	\end{equation*}
	and, as a consequence,
	\[
	C_\infty:=\lim\limits_{t\to\infty} C(t) = \frac{2^{d-2}\mathcal{M}_{\mathcal{R}}}{\sqrt{\pi} }.
	\]
\end{lem}

\begin{proof}
	For the sake of simplicity, we will consider $t(\tau)=e^\tau$ in the computations instead of the real change of variables $t(\tau)=e^\tau-1$. This does not affect the results, it is just a translation in time; instead of $t\in [0,\infty)$ we will work on the time interval $t\in[1, \infty)$, and then a map $t\to t+1$ will undo de translation.
	
	In order to estimate the function $C(t)$ we notice that $V$ must have constant mass, since a solution of equation~\eqref{eq:main_spherical_coordinates_symmetric_datum} does so. Therefore,
	\begin{equation}\label{eq:equality_C_radial_solutions}
		\begin{aligned}
			\mathcal{M}_{\mathcal{R}} = \int_{\rho_0(\tau)}^\infty V(\tau, \rho)\ \d \rho &= \frac{C(t)}{\sqrt{t}}\int_{0}^\infty \sinh^{d-1}(r)\cdot e^{- \frac{\left(r + (d-1)t\right)^2}{4t}}\ \d r
		\end{aligned}
	\end{equation}	
	for all $t\geq 1$. Using implicit derivation in the equation in~\eqref{eq:equality_C_radial_solutions} it is possible to estimate $\partial_t C(t)/C(t)$. It can be verified, after some computations, that
	\begin{equation}\label{eq:quotient_C(t)}
		\frac{\partial_t C(t)}{C(t)} = \frac{(d-1)^2t+2}{4t} - \frac{1}{4t^2}\cdot \frac{\mathcal{I}}{\mathcal{II}},
	\end{equation}
	where
	\[
	\mathcal{I}:= \int_0^\infty r^2 (1-e^{-2r})^{d-1}\ e^{- \frac{\left(r - (d-1)t\right)^2}{4t}}\ \d r,\quad  \mathcal{II}:= \int_0^\infty (1-e^{-2r})^{d-1}\ e^{- \frac{\left(r - (d-1)t\right)^2}{4t}}\ \d r.
	\]
	
	Let us begin with a lower estimate  for ~\eqref{eq:quotient_C(t)}.
	%	
	%	First, we estimate the quotient~\eqref{eq:quotient_C(t)} on a compact interval $[1, T^*_d]$, for a $T^*_d\geq 1$ to be chosen later. Since the quotient is finite at time $t=1$ and the functions of time involved in its formulas are continuous , it is clear that there must exist a constant $K(T^*_d)$ such that
	%	\[
	%	-K(T^*_d)\leq \frac{\partial_t C(t)}{C(t)}\leq  K(T^*_d) \quad\text{for all}\quad t\in [1, T^*_d].
	%	\]
	%	
	%	Now, for times bigger than $T^*_d$.
	We estimate each integral separately. First, with a change of variables and using that $1-e^{-2r}<1$ ,
	\[
	\begin{aligned}
		\mathcal{I} &\leq  2\sqrt{t}\int_{-\sqrt{\lambda_1 t}}^\infty \left(2\sqrt{t}z + (d-1)t\right)^2 \ e^{-z^2}\ \d z\\[10pt]
		& = \underbrace{8t^\frac{3}{2}\int_{-\sqrt{\lambda_1 t}}^\infty z^2 \ e^{-z^2}\ \d z}_{I_1} + \underbrace{8(d-1)t^2 \int_{-\sqrt{\lambda_1 t}}^\infty z \ e^{-z^2}\ \d z}_{I_2} + \underbrace{2 (d-1)^2 t^\frac{5}{2} \int_{-\sqrt{\lambda_1 t}}^\infty  e^{-z^2}\ \d z}_{I_3},
	\end{aligned}
	\]
	where $\lambda_1$ comes from~\eqref{eq:definition_bottom_spectrum_laplacian}. The integral $I_1$ can be bounded with an integration by parts and using that
	\[
	\int_{-\sqrt{\lambda_1 t}}^\infty e^{-z^2}\d z\leq  \int_{-\infty}^\infty e^{-z^2}\d z = \sqrt{\pi}
	\]
	by
	\[
	I_1\leq 2\sqrt{t}\left(2\sqrt{\pi}\ t - (d-1)t^\frac{3}{2}e^{-\lambda_1 t}  \right),
	\]
	while $I_2$ can be computed explicitly as
	\[
	I_2 = 4(d-1)t^2 e^{-\lambda_1 t}
	\]
	and
	\[
	I_3 \leq 2 (d-1)^2 \sqrt{\pi}\  t^\frac{5}{2}.
	\]
	In total,
	\[
	\mathcal{I}\leq 2\sqrt{t}\left((d-1)^2\sqrt{\pi}\ t^2  +2\sqrt{\pi}\ t  + (d-1)t^\frac{3}{2}e^{-\lambda_1 t} \right)
	\]
	Second, we compute for $\mathcal{II}$,
	\[
	\begin{aligned}
		\mathcal{II} &= 2\sqrt{t} \int_{-\sqrt{\lambda_1 t}}^{\infty} \left(1- e^{-2(2\sqrt{t}z+(d-1)t)}\right)^{d-1} e^{-z^2}\ \d z\\[10pt]
		&\geq 2\sqrt{t} \int_{-\frac{\sqrt{\lambda_1 t}}{2}}^{\infty} \left(1- e^{-2(2\sqrt{t}z+(d-1)t)}\right)^{d-1} e^{-z^2}\ \d z\\[10pt]
		& \geq  2\sqrt{t} \left(1- e^{-(d-1)t}\right)^{d-1} \int_{-\frac{\sqrt{\lambda_1 t}}{2}}^{\infty}  e^{-z^2}\ \d z\\[10pt]
		&= 2\sqrt{t} \left(1- e^{-(d-1)t}\right)^{d-1}\left[ \int_{-\infty}^{\infty}  e^{-z^2}\ \d z - \int_{-\infty}^{-\frac{\sqrt{\lambda_1 t}}{2}}  e^{-z^2}\ \d z \right]\\[10pt]
		&= 2\sqrt{t} \left(1- e^{-(d-1)t}\right)^{d-1}\left[\sqrt{\pi} - \int_{-\infty}^{-\frac{\sqrt{\lambda_1 t}}{2}}  e^{-z^2}\ \d z \right]\\[10pt]
		&\geq 2\sqrt{t} \left(1- e^{-(d-1)t}\right)^{d-1}\left[\sqrt{\pi} +\frac{1}{\sqrt{\lambda_1 t}} \int_{-\infty}^{-\frac{\sqrt{\lambda_1 t}}{2}} 2z e^{-z^2}\ \d z \right]\\[10pt]
		&= 2\sqrt{t} \left(1- e^{-(d-1)t}\right)^{d-1}\left[\sqrt{\pi} - \frac{1}{\sqrt{\lambda_1 t}}e^{-{\frac{\lambda_1 t}{4}}} \right].
	\end{aligned}
	\]
	In total, we arrive to
	\[
	\begin{aligned}
		\frac{\partial_t C(t)}{C(t)}&\geq \frac{(d-1)^2t+2}{4t} - \frac{(d-1)^2\sqrt{\pi}\ t^2  +2\sqrt{\pi}\ t  + (d-1)t^\frac{3}{2}e^{-\lambda_1 t}}{4t^2\left(1-e^{-(d-1)t}\right)^{d-1} \left(\sqrt{\pi} - \frac{1}{\sqrt{\lambda_1 t}}e^{-\frac{\lambda_1 t}{4}}\right)}\\[10pt]
		&=\frac{(d-1)^2t+2}{4t}\cdot h(t) - \frac{\left(\frac{2\sqrt{t}}{\sqrt{\lambda_1}}+2(d-1)t^{\frac{3}{2}}\right)e^{-\frac{\lambda_1 t}{4}} + (d-1)t^\frac{3}{2}e^{-\lambda_1 t} }{4t^2\left(1-e^{-(d-1)t}\right)^{d-1} \left(\sqrt{\pi} - \frac{1}{\sqrt{\lambda_1 t}}e^{-\frac{\lambda_1 t}{4}}\right)},
	\end{aligned}
	\]
	where
	\begin{equation}\label{eq:def_h(t)}
		h(t):= \left(1-\left(1-e^{-(d-1)t}\right)^{-(d-1)}\right)<0.
	\end{equation}
	Since
	\[
	\lim\limits_{t\to\infty} \frac{h(t)}{-(d-1)e^{-(d-1)t}} = 1,
	\]
	it is left as an exercise for the reader to check that 
	\[
	h(t)\geq -2(d-1)e^{-(d-1)t}\quad\text{for all}\quad t\geq 1.
	\]
	Therefore, for $t\geq 1$,
	\[
	\frac{\partial_t C(t)}{C(t)}\geq -\frac{(d-1)^3t+2(d-1)}{2t} e^{-(d-1)t} - \frac{\left(\frac{2\sqrt{t}}{\sqrt{\lambda_1}}+2(d-1)t^{\frac{3}{2}}\right)e^{-\frac{\lambda_1 t}{4}} + (d-1)t^\frac{3}{2}e^{-\lambda_1 t} }{4t^2\left(1-e^{-(d-1)t}\right)^{d-1} \left(\sqrt{\pi} - \frac{1}{\sqrt{\lambda_1 t}}e^{-\frac{\lambda_1 t}{4}}\right)}.
	\]
	Thus, if we define 
	\begin{equation}\label{eq:def_m_d}
		m_d := \min \left(d-1, \frac{(d-1)^2}{16}\right),
	\end{equation}
	then there must exist a positive constant $k_d$ such that
	\[
	\frac{\partial_t C(t)}{C(t)}\geq -k_d\ e^{-m_d t} \quad\text{for all}\quad t\geq 1.
	\]
	
	Now, for the upper bound  for ~\eqref{eq:quotient_C(t)}. For the sake of brevity, we will omit computations similar to the ones in the previous step. We estimate again each integral separately. First, by similar arguments as before,
	\[
	\mathcal{I}\geq 2\sqrt{t}\left(1-e^{-(d-1)t}\right)^{d-1}\int_{-\frac{\sqrt{\lambda_1 t}}{2}}^\infty \left(4tz^2 +(d-1)^2t^2  +4(d-1)t^\frac{3}{2}  \right) e^{-z^2}\ \d z.
	\]
	We use now the following bound, left for the reader as an exercise. For every $a>0$,
	\[
	\int_{-a}^\infty e^{-z^2}\ \d z = \int_{-\infty}^\infty e^{-z^2}\ \d z - \int_{-\infty}^{-a} e^{-z^2}\ \d z\geq \sqrt{\pi} - \frac{1}{2a}e^{-a^2}.
	\]
	Using this, and that $\sqrt{\lambda_1}=(d-1)/2$ by definition, it is not hard to arrive to
	\[
	\mathcal{I}\geq 2\sqrt{t}\left(1-e^{-(d-1)t}  \right)^{d-1}\left[\sqrt{\pi}(d-1)^2t^2+2\sqrt{\pi}t + e^{-\frac{\lambda_1 t}{4}}\left(   \frac{d-1}{2}t^\frac{3}{2} - \frac{2}{\sqrt{\lambda_1}}\sqrt{t}  \right)  \right].
	\]
	For the integral $\mathcal{II}$ we simply compute
	\[
	\mathcal{II}\leq 2\sqrt{t}\int_{-\sqrt{\lambda_1 t}}^\infty e^{-z^2}\ \d z\leq 2\sqrt{t}\int_{-\infty}^\infty e^{-z^2}\ \d z = 2\sqrt{\pi t}.
	\]
	Substituting in~\eqref{eq:quotient_C(t)} we arrive, after some computations, to
	\[
	\frac{\partial_t C(t)}{C(t)}\leq \left( \frac{(d-1)^2}{4} + \frac{1}{2t} \right)h(t) + \left(1-e^{-(d-1)t}  \right)^{d-1} \left( \frac{d-1}{2}t^\frac{3}{2} - \frac{2}{\sqrt{\lambda_1}}\sqrt{t} \right) \frac{e^{-\frac{\lambda_1 t}{4}}}{4t^2\sqrt{\pi}},
	\]
	where $h(t)$ is defined as in~\eqref{eq:def_h(t)}. From here it follows that
	\[
	\frac{\partial_t C(t)}{C(t)}\leq \left( \frac{(d-1)^3}{2} + \frac{1}{t} \right)e^{-(d-1)t} + \left( \frac{d-1}{2}t^\frac{3}{2} - \frac{2}{\sqrt{\lambda_1}}\sqrt{t} \right) \frac{e^{-\frac{\lambda_1 t}{4}}}{4t^2\sqrt{\pi}},
	\]
	and from here, with $m_d$ defined as in~\eqref{eq:def_m_d},
	\[
	\frac{\partial_t C(t)}{C(t)}\leq k_d e^{-m_d t}\quad\text{for all}\quad t>1.
	\]
	Putting all together we obtain~\eqref{eq:bound_quotient_C(t)}.
	
	Along the previous steps, we have proven that
	\[
	2\sqrt{t} \left(1- e^{-(d-1)t}\right)^{d-1}\left[\sqrt{\pi} - \frac{1}{\sqrt{\lambda_1 t}}e^{-{\frac{\lambda_1 t}{4}}} \right]\leq  \mathcal{II}\leq  2\sqrt{\pi t}
	\]
	With this bounds and equation~\eqref{eq:equality_C_radial_solutions} it is easy to conclude~\eqref{eq:limit_C(t)}.
\end{proof}

\end{document}
